% Part: second-order-logic
% Chapter: metatheory
% Section: undecidability-and-axiomatizability

\documentclass[../../../include/open-logic-section]{subfiles}

\begin{document}

\olfileid{sol}{met}{nax}

\olsection{ద్వితీయ-స్థాయి తర్కం స్వీకృతీకరించదగినది కాదు}


\begin{thm}
\ollabel{thm:sol-undecidable}
ద్వితీయ-స్థాయి తర్కం అనిర్ణయనీయమైనది.
\end{thm}

\begin{proof}
ఒక మొదటిస్థాయి \tetoken{వాక్యం}{!!{sentence}} మొదటిస్థాయి తర్కంలో
చెల్లుబాటైనప్పుడు, అప్పుడే మాత్రమే అది ద్వితీయ-స్థాయి తర్కంలోనూ చెల్లుబాటవుతుంది.
మొదటిస్థాయి తర్కం అనిర్ణయనీయమైనది.
\end{proof}

\begin{thm}
\ollabel{cor:sol-not-axiomatizable}
ద్వితీయ-స్థాయి తర్కానికి నిర్దుష్టమైన, సంపూర్ణమైన \tetoken{వ్యుత్పత్తి}{!!{derivation}}
వ్యవస్థ ఏదీ లేదు.
\end{thm}

\begin{proof}
$!A$ అంకగణిత భాషలోని \tetoken{ఒక వాక్యం}{!!a{sentence}} అనుకుందాం.
$\Sat{N}{!A}$ నెరవేరినప్పుడు, అప్పుడే మాత్రమే $\Th{PA^2} \Entails !A$ నెరవేరుతుంది. $\Th{PA^2}$లోని
తొమ్మిది స్వీకృతాల సంయోగాన్ని $!P$గా తీసుకుందాం. $\Th{PA^2} \Entails !A$
నెరవేరినప్పుడు, అప్పుడే మాత్రమే $\Entails !P \lif !A$ నెరవేరుతుంది. అంటే ప్రతి అంకగణిత నిర్మాణంలోనూ
$\Sat{M}{!P \lif !A}$.
మూల వాక్యంలోని బద్ధ చరాల పేర్లను అవసరమైతే ముందుగా మార్చుదాం.
అప్పుడు కింద ప్రవేశపెట్టే చరాలేవీ ఆ వాక్యంలో ఉండవు.
ఇప్పుడు $\Obj{0}$ స్థానంలో $z$ను, $\prime$ స్థానంలో ఏకస్థాన ప్రమేయ
చరం $u$ను, $+$, $\times$ల స్థానంలో వరుసగా ద్విస్థాన ప్రమేయ చరాలు
$u'$, $u''$లను, $<$ స్థానంలో ద్విస్థాన సంబంధ చరం~$L$ను ప్రతిస్థాపించి,
సార్వత్రికంగా పరిమాణీకరించడం వల్ల వచ్చే
\tetoken{వాక్యాన్ని}{!!{sentence}}
$\lforall[z][\lforall[u][\lforall[u'][\lforall[u''][\lforall[L][(!P' \lif !A')]]]]]$
పరిశీలించండి. అసలు వాక్యం చెల్లుబాటైనప్పుడు, అప్పుడే మాత్రమే ఇది శుద్ధ
ద్వితీయ-స్థాయి తర్కంలో చెల్లుబాటవుతుంది. అసలు వాక్యం చెల్లుబాటయ్యేది
$\Th{PA^2} \Entails !A$ నెరవేరినప్పుడే. ఇది $\Sat{N}{!A}$ నెరవేరడంతో సమానం. కాబట్టి ద్వితీయ-స్థాయి తర్కానికి నిర్దుష్టమైన,
సంపూర్ణమైన నిరూపణ వ్యవస్థ ఉంటే,~$\Struct{N}$లో సత్యమైన
\tetoken{వాక్యాల}{!!{sentence}s} గణనీయ జాబితాను ఇచ్చే
$f\colon \Nat \to \Sent[L_A]$ను నిర్వచించడానికి దాన్ని ఉపయోగించవచ్చు.
ఈ ప్రమేయానికి~$\Th{Q}$లో ఏదో మొదటిస్థాయి సూత్రం~$!B_f(x, y)$
ప్రాతినిధ్యం వహిస్తుంది. అప్పుడు \tetoken{సూత్రం}{!!{formula}}
$\lexists[x][!B_f(x, y)]$ అనేది $\Struct{N}$లో సత్యమైన మొదటిస్థాయి
\tetoken{వాక్యాల}{!!{sentence}s} సమితిని నిర్వచిస్తుంది. ఇది టార్స్కీ సిద్ధాంతానికి
విరుద్ధం.
\end{proof}



\end{document}
